\chapter{Systems of Linear Equations and Their Geometry}

A \emph{system of linear equations} is the single most ordinary object in
all of mathematics. You have met one before, probably in middle school,
disguised as a word problem: ``two pencils and a notebook cost \$5, three
pencils and two notebooks cost \$8, find the price of each.'' You may have
solved it by substitution or by the comforting ritual of ``add the two
equations to cancel a variable.'' It worked, and you moved on.

This chapter asks you to slow down and \emph{look} at what you were doing,
because hidden inside that humble word problem is the entire skeleton of
linear algebra. A system like
\[
  \begin{aligned}
    2x + 1y &= 5,\\
    3x + 2y &= 8
  \end{aligned}
\]
can be \emph{seen} in two different ways, and the whole subject grows out of
holding both pictures in your head at once. One picture is about \emph{rows}
-- lines and planes crossing. The other is about \emph{columns} -- arrows
adding up. They describe the same equations and give the same answer, yet
they feel like two different universes. Learning to flip between them at will
is the skill that makes everything after this chapter click.

We will draw both pictures carefully, in two and three dimensions, with
small numbers you can check by hand. Then we will ask the three questions
that organize all of linear algebra:
\begin{itemize}
  \item When does a system have \emph{exactly one} solution?
  \item When does it have \emph{none}?
  \item When does it have \emph{infinitely many}?
\end{itemize}
The astonishing fact -- the one this chapter exists to convince you of -- is
that those are the \emph{only} three possibilities, and that each of them
has a clean geometric meaning. The words \emph{singular} and
\emph{nonsingular} will turn out to be names for ``the columns collapsed''
versus ``the columns stayed spread out.'' Hold on to that phrase; we will
earn it.

\section{Two pictures of one system}

Let us fix a concrete system and never let go of it:
\begin{equation}
  \label{ch06:flagship}
  \begin{aligned}
    2x + 1y &= 5,\\
    3x + 2y &= 8.
  \end{aligned}
\end{equation}
There are two ways to read these symbols, and both are correct and useful.
Most people are taught only the first and never told the second exists, which
is a shame, because the second is the one that grows up into the rest of the
book.

\subsection{The row picture: lines that cross}

Read \eqref{ch06:flagship} one row at a time. The first row, $2x+1y=5$, is a
\emph{constraint} on the pair $(x,y)$. The set of all points $(x,y)$ that
satisfy it is a straight line in the plane. The second row, $3x+2y=8$, is a
different constraint, a different line. A pair $(x,y)$ solves the
\emph{system} exactly when it sits on \emph{both} lines at once -- that is,
at the point where the two lines cross.

\begin{intuitionbox}
The row picture turns ``solve the system'' into ``find where the lines
meet.'' Each equation is a line (in $2$ unknowns) or a plane (in $3$
unknowns). A solution is a point common to all of them. Geometry has quietly
replaced algebra.
\end{intuitionbox}

\begin{figure}[H]
\centering
\begin{tikzpicture}[scale=0.95]
  \draw[->,thick] (-0.5,0) -- (4.5,0) node[right] {$x$};
  \draw[->,thick] (0,-0.5) -- (0,4.7) node[above] {$y$};
  \foreach \i in {1,2,3,4} {\draw[gray!40] (\i,-0.08)--(\i,0.08); \draw[gray!40] (-0.08,\i)--(0.08,\i);}
  % line 1: 2x + y = 5  -> y = 5 - 2x  passes (0,5),(2.5,0)
  \draw[vecblue,very thick,domain=0.3:2.55] plot (\x,{5-2*\x});
  \node[vecblue,above right] at (0.5,4.0) {\small $2x+y=5$};
  % line 2: 3x + 2y = 8 -> y = (8-3x)/2 passes (0,4),(8/3,0)
  \draw[vecred,very thick,domain=-0.1:2.65] plot (\x,{(8-3*\x)/2}) node[below right,vecred] {\small $3x+2y=8$};
  % intersection x=2, y=1
  \filldraw[black] (2,1) circle (2.4pt) node[above right] {\small $(2,1)$};
\end{tikzpicture}
\caption{The row picture of system~\eqref{ch06:flagship}. Each equation is a
line; the solution is the single point $(2,1)$ where they cross.}
\label{ch06:rowpic}
\end{figure}

You can read the answer right off Figure~\ref{ch06:rowpic}: the lines cross
at $(x,y)=(2,1)$. Check it: $2(2)+1=5$, yes; $3(2)+2(1)=8$, yes. Two lines in
a plane, if they are not parallel, meet in exactly one point. That ``if they
are not parallel'' is the seed of everything that can go wrong, and we will
return to it.

\subsection{The column picture: arrows that add up}

Now read the \emph{same} system \eqref{ch06:flagship} a different way -- not
row by row, but \emph{column by column}. Stack the two equations
into a single vector equation:
\[
  \colvec{2x+1y\\3x+2y} = \colvec{5\\8}.
\]
The left side is a sum of two pieces, one carrying $x$ and one carrying $y$:
\[
  x\colvec{2\\3} + y\colvec{1\\2} = \colvec{5\\8}.
\]
Look hard at this. It says: take the vector $\vec{a}_1=\colvec{2\\3}$, the
vector $\vec{a}_2=\colvec{1\\2}$, and the target $\vec{b}=\colvec{5\\8}$. The
question is no longer ``where do two lines cross?'' It is now:

\begin{focusbox}
\emph{How much of $\vec{a}_1$ and how much of $\vec{a}_2$ must I add to land
exactly on $\vec{b}$?} The unknowns $x$ and $y$ are the \emph{amounts}. A
linear system asks you to build the target vector $\vec b$ out of the
\emph{columns} of the matrix.
\end{focusbox}

The columns of the system are the arrows $\vec a_1,\vec a_2$. We are looking
for the recipe $x\,\vec a_1 + y\,\vec a_2$ that reaches $\vec b$. This is a
\emph{linear combination} of the columns -- exactly the operation you met
when we first talked about span. Solving a linear system \emph{is} asking
whether $\vec b$ lies in the span of the columns, and if so, with what
weights.

\begin{figure}[H]
\centering
\begin{tikzpicture}[scale=0.68,>=Stealth]
  \draw[help lines,gray!20] (-0.5,-0.5) grid (6.5,8.8);
  \draw[->,gray!60] (-0.5,0) -- (6.8,0) node[right] {$x$};
  \draw[->,gray!60] (0,-0.5) -- (0,9.0) node[above] {$y$};
  % columns
  \draw[->,vecblue,very thick] (0,0) -- (2,3) node[above left=1pt] {\small $\vec a_1=\colvec{2\\3}$};
  \draw[->,targetgreen!80!black,very thick] (0,0) -- (1,2) node[below right=1pt] {\small $\vec a_2=\colvec{1\\2}$};
  % scaled: 2 a1 = (4,6); then + 1 a2 = (5,8)
  \draw[->,vecblue,thick,dashed] (0,0) -- (4,6) node[midway,left=1pt] {\small $2\vec a_1$};
  \draw[->,targetgreen!80!black,thick,dashed] (4,6) -- (5,8) node[midway,right=1pt] {\small $1\vec a_2$};
  \draw[->,vecred,ultra thick] (0,0) -- (5,8) node[midway,above left=2pt] {\small $\vec b=\colvec{5\\8}$};
  \filldraw[vecred] (5,8) circle (2.8pt) node[above right=1pt,vecred] {\small $(5,8)$};
  \filldraw[black] (0,0) circle (1.8pt);
\end{tikzpicture}
\caption{The column picture of system~\eqref{ch06:flagship}. We need
$x\,\vec a_1 + y\,\vec a_2 = \vec b$. The answer $x=2,\ y=1$ says: walk twice
along $\vec a_1$ (reaching $(4,6)$), then once along $\vec a_2$, and you arrive exactly at
$\vec b=\colvec{5\\8}$.}
\label{ch06:colpic}
\end{figure}

The two pictures use the very same numbers, but they organize them
orthogonally (pun intended). In the row picture, the rows $(2,1)$ and
$(3,2)$ are the stars and they appear as \emph{lines}. In the column
picture, the columns $(2,3)$ and $(1,2)$ are the stars and they appear as
\emph{arrows}. The right-hand side $(5,8)$ is a \emph{pair of heights} (one
per line) in the row picture, but a single \emph{target arrow} in the column
picture. Same data, transposed in your imagination.

\begin{trapbox}
The single most common confusion in week one of linear algebra: mixing up
the rows and the columns of $\colvec{2 & 1\\3 & 2}$. The \emph{rows}
$(2,1),(3,2)$ are coefficients of the two equations -- they make the
\emph{lines}. The \emph{columns} $\colvec{2\\3},\colvec{1\\2}$ are the
arrows you combine -- they make the \emph{span}. When you read off a column
of a matrix, read \emph{down}.
\end{trapbox}

\subsection{Matrix form: the whole thing as $A\vec x = \vec b$}

Both pictures collapse into one tidy statement. Collect the columns
side by side into a \emph{matrix} $A$, stack the unknowns into a vector
$\vec x$, and the targets into $\vec b$:
\[
  A=\colvec{2 & 1\\3 & 2},\qquad
  \vec x=\colvec{x\\y},\qquad
  \vec b=\colvec{5\\8},
\]
so that the system is simply
\[
  A\vec x = \vec b.
\]
The notation hides the duality but does not destroy it. The product
$A\vec x$ can be computed two ways, and they are exactly our two pictures:
\begin{itemize}
  \item \textbf{Row way:} the $i$-th entry of $A\vec x$ is row $i$ of $A$
  dotted with $\vec x$. Setting each dotted result equal to the matching
  entry of $\vec b$ gives back the original equations (the lines).
  \item \textbf{Column way:} $A\vec x = x\,(\text{col }1) + y\,(\text{col
  }2)$, a linear combination of the columns (the arrows).
\end{itemize}
Keep both readings of $A\vec x$ alive. When we discuss whether a solution
exists, the column reading dominates; when we actually \emph{compute} one by
elimination, the row reading dominates.

\begin{strangbox}
\emph{Before Strang:} when you open his book, the very first lecture is
``the geometry of linear equations,'' and it is built on precisely this
row-versus-column split. Strang famously says the column picture is the one
students under-appreciate, and that ``$A\vec x$ is a combination of the
columns of $A$'' is the sentence to tattoo on your arm. If you internalize
this chapter, his Lecture~1 will feel like a friend you already know.
\emph{After Strang:} the column picture is also where the four fundamental
subspaces are born -- the set of all reachable $\vec b$ is the
\emph{column space}, which you will meet soon.
\end{strangbox}

\section{A fully worked $2\times 2$, both ways}

Let us solve \eqref{ch06:flagship} from scratch with bare algebra, so the
geometry has numbers under it.

\begin{example}[Substitution, the row way]
\label{ch06:ex-sub}
From the first equation, $y = 5 - 2x$. Substitute into the second:
\[
  3x + 2(5-2x) = 8 \;\Longrightarrow\; 3x + 10 - 4x = 8 \;\Longrightarrow\;
  -x = -2 \;\Longrightarrow\; x = 2.
\]
Then $y = 5 - 2(2) = 1$. So $(x,y)=(2,1)$, matching the crossing point in
Figure~\ref{ch06:rowpic}. Substitution is the row picture made mechanical:
we used equation one to pin $y$ to $x$ (slide along the first line), then
forced the second line.
\end{example}

\begin{example}[Combination, the column way]
\label{ch06:ex-col}
We want weights $x,y$ with $x\colvec{2\\3}+y\colvec{1\\2}=\colvec{5\\8}$.
Eliminate the second column from the top: multiply the system's bottom entry
by something and subtract. Concretely, the inverse recipe (we will derive it
below) gives
\[
  x = \frac{2\cdot 5 - 1\cdot 8}{2\cdot 2 - 1\cdot 3}=\frac{10-8}{4-3}=2,
  \qquad
  y = \frac{2\cdot 8 - 3\cdot 5}{2\cdot 2 - 1\cdot 3}=\frac{16-15}{1}=1.
\]
Same answer, $x=2,\ y=1$. Notice the denominator $2\cdot2-1\cdot3=1$. That
number -- the \emph{determinant} of $A$ -- is about to become the hero of the
chapter. When it is nonzero, the columns are spread out and a unique recipe
exists. When it is zero, they have collapsed onto one line, and trouble
begins.
\end{example}

\begin{labbox}
Solve $\colvec{1 & 2\\1 & 3}\colvec{x\\y}=\colvec{4\\5}$ \textbf{twice}:
once by substitution (find the crossing point of the two lines) and once by
combination (what multiples of $\colvec{1\\1}$ and $\colvec{2\\3}$ reach
$\colvec{4\\5}$?). You should get $x=2,\ y=1$ both ways. Then sketch both
pictures on scrap paper. The point is not the answer; it is feeling the two
machines grind to the same place.
\end{labbox}

\section{The three possibilities, and why there are only three}

So far our lines crossed once and our columns reached the target once. Life
is not always so kind. Change a few numbers and the system can have
\emph{no} solution or \emph{infinitely many}. The miracle is that those are
the only alternatives. Let us see all three in the cleanest possible
setting, $2\times 2$, and watch both pictures react.

\subsection{One solution: lines cross, columns spread}

This is the case we just did. Two lines with \emph{different slopes} cross
at exactly one point. In the column picture, the two arrows
$\vec a_1,\vec a_2$ point in \emph{different directions}, so by scaling and
adding them you can reach \emph{any} target $\vec b$ in the plane -- and in
exactly one way.

\begin{figure}[H]
\centering
\begin{tikzpicture}[scale=0.8]
  \begin{scope}
    \node at (1.5,3.6) {\small Row: lines cross};
    \draw[->,gray] (-0.4,0)--(3.4,0); \draw[->,gray] (0,-0.4)--(0,3.4);
    \draw[vecblue,very thick] (-0.2,2.7)--(3.2,0.15);
    \draw[vecred,very thick] (-0.2,0.4)--(3.0,3.2);
    \filldraw (1.55,1.55) circle (2pt);
  \end{scope}
  \begin{scope}[xshift=6.2cm]
    \node at (1.5,3.6) {\small Column: arrows spread};
    \draw[->,gray] (-0.4,0)--(3.4,0); \draw[->,gray] (0,-0.4)--(0,3.4);
    \draw[->,vecblue,very thick] (0,0)--(2.4,0.7) node[right]{\small $\vec a_1$};
    \draw[->,targetgreen!80!black,very thick] (0,0)--(0.7,2.4) node[above]{\small $\vec a_2$};
  \end{scope}
\end{tikzpicture}
\caption{\emph{Nonsingular} case. Different slopes / different directions.
Exactly one solution for every $\vec b$.}
\label{ch06:one}
\end{figure}

\subsection{No solution: parallel lines, collinear columns missing the target}

Take a system whose two lines are \emph{parallel but not equal}:
\begin{equation}
  \label{ch06:nosol}
  \begin{aligned}
    1x + 2y &= 3,\\
    1x + 2y &= 7.
  \end{aligned}
\end{equation}
The rows have the same left side, $x+2y$, but demand two different values,
$3$ and $7$. No point can make $x+2y$ equal $3$ and $7$ at once. The two
lines never meet -- they are parallel. \emph{No solution.}

Now flip to columns. The system is
$x\colvec{1\\1}+y\colvec{2\\2}=\colvec{3\\7}$. But $\colvec{1\\1}$ and
$\colvec{2\\2}$ point the \emph{same way} -- the second is just twice the
first. Every combination $x\colvec{1\\1}+y\colvec{2\\2}=(x+2y)\colvec{1\\1}$
lands somewhere on the single line through $\colvec{1\\1}$. But the target
$\colvec{3\\7}$ is \emph{off} that line (it would need its two entries equal,
and $3\neq7$). So no combination reaches it.

\begin{figure}[H]
\centering
\begin{tikzpicture}[scale=0.8]
  \begin{scope}
    \node at (1.5,3.6) {\small Row: parallel lines};
    \draw[->,gray] (-0.4,0)--(3.4,0); \draw[->,gray] (0,-0.4)--(0,3.4);
    \draw[vecblue,very thick] (-0.3,1.65)--(3.3,-0.15);
    \draw[vecred,very thick] (-0.3,3.05)--(3.3,1.25);
  \end{scope}
  \begin{scope}[xshift=6.2cm]
    \node at (1.6,3.6) {\small Column: collinear, $\vec b$ off the line};
    \draw[->,gray] (-0.4,0)--(3.4,0); \draw[->,gray] (0,-0.4)--(0,3.4);
    \draw[gray!40,dashed] (0,0)--(3.2,3.2);
    \draw[->,vecblue,very thick] (0,0)--(1.0,1.0) node[below right]{\small $\vec a_1$};
    \draw[->,targetgreen!80!black,very thick] (0,0)--(2.0,2.0) node[below right]{\small $\vec a_2$};
    \draw[->,vecred,ultra thick] (0,0)--(1.1,2.7) node[above]{\small $\vec b$};
  \end{scope}
\end{tikzpicture}
\caption{\emph{No solution.} The lines are parallel; equivalently, the
columns are collinear (both on the dashed line) and the target $\vec b$ does
not lie on that line.}
\label{ch06:none}
\end{figure}

\subsection{Infinitely many: same line twice, target on the column line}

Keep the columns collinear, but now put the target \emph{on} their line:
\begin{equation}
  \label{ch06:infsol}
  \begin{aligned}
    1x + 2y &= 3,\\
    2x + 4y &= 6.
  \end{aligned}
\end{equation}
The second equation is just twice the first -- it is the \emph{same line}
written twice. Any point on the line $x+2y=3$ satisfies both. There are
infinitely many such points: $x=3-2t,\ y=t$ for every real $t$. In the
column picture, $\colvec{1\\2}$ and $\colvec{2\\4}$ are collinear (the second
is twice the first), and the target $\colvec{3\\6}$ \emph{is} on that line,
so it can be reached -- but in infinitely many ways, because you can trade
amount of one column for amount of the other.

\begin{figure}[H]
\centering
\begin{tikzpicture}[scale=0.8]
  \begin{scope}
    \node at (1.5,3.6) {\small Row: one line, twice};
    \draw[->,gray] (-0.4,0)--(3.4,0); \draw[->,gray] (0,-0.4)--(0,3.4);
    \draw[structurepurple,very thick] (-0.3,3.3)--(3.3,-0.3);
    \filldraw (0.6,2.4) circle (1.5pt);
    \filldraw (1.6,1.4) circle (1.5pt);
    \filldraw (2.6,0.4) circle (1.5pt);
  \end{scope}
  \begin{scope}[xshift=6.2cm]
    \node at (1.6,3.6) {\small Column: $\vec b$ on the line};
    \draw[->,gray] (-0.4,0)--(3.4,0); \draw[->,gray] (0,-0.4)--(0,3.4);
    \draw[gray!40,dashed] (0,0)--(1.6,3.2);
    \draw[->,vecblue,very thick] (0,0)--(0.7,1.4) node[right]{\small $\vec a_1$};
    \draw[->,targetgreen!80!black,very thick] (0,0)--(1.4,2.8) node[right]{\small $\vec a_2$};
    \draw[->,vecred,ultra thick] (0,0)--(1.05,2.1) node[left]{\small $\vec b$};
  \end{scope}
\end{tikzpicture}
\caption{\emph{Infinitely many solutions.} The two equations are the same
line; the columns are collinear and $\vec b$ lies on their common line, so it
is reachable in infinitely many ways.}
\label{ch06:inf}
\end{figure}

\begin{focusbox}
Read off the pattern. Whether the system has $0$, $1$, or $\infty$ solutions
is decided by two yes/no questions:
\begin{enumerate}
  \item \textbf{Are the columns spread out} (point in different
  directions)? If \emph{yes}, you can reach \emph{every} $\vec b$, uniquely:
  one solution, always.
  \item If the columns have \textbf{collapsed} (collinear), then it depends
  on $\vec b$: if $\vec b$ lies on their line, \emph{infinitely many}
  solutions; if not, \emph{none}.
\end{enumerate}
``Spread out vs.\ collapsed'' is the heart of the matter. It has a name:
\emph{nonsingular} vs.\ \emph{singular}.
\end{focusbox}

\subsection{Why never exactly two?}

A student always asks: why can't a system have exactly two solutions, or
seventeen? Here is the clean reason, and it is worth seeing once.

\begin{proposition}[Zero, one, or infinity -- never in between]
\label{ch06:zoi}
If a linear system $A\vec x=\vec b$ has two distinct solutions, it has
infinitely many.
\end{proposition}

\begin{proof}
Suppose $\vec p$ and $\vec q$ are both solutions, $\vec p\neq\vec q$, so
$A\vec p=\vec b$ and $A\vec q=\vec b$. Consider the whole \emph{line}
through them: for any real number $t$, set
$\vec x_t=\vec p + t(\vec q-\vec p)$. Then, using that $A$ distributes over
sums and pulls out scalars (linearity),
\[
  A\vec x_t = A\vec p + t\,A(\vec q-\vec p)
            = \vec b + t(A\vec q - A\vec p)
            = \vec b + t(\vec b-\vec b) = \vec b.
\]
So \emph{every} point $\vec x_t$ on the line is a solution, and since
$\vec q-\vec p\neq\vec 0$ these are all different. Two solutions force a
whole line of them -- hence infinitely many.
\end{proof}

\begin{intuitionbox}
Geometrically: the solution set of a linear system is never a scattering of
isolated dots. It is always ``flat'' -- a single point, or a whole line, or
a plane, or empty. Linearity forbids curved or finite-but-multiple solution
sets. That flatness is exactly why the answer count is $0$, $1$, or
$\infty$ and nothing else.
\end{intuitionbox}

\section{Singular versus nonsingular: the determinant in 2D}

We keep saying ``spread out'' and ``collapsed.'' Let us make it a number.
For a $2\times 2$ matrix
\[
  A=\colvec{a & b\\ c & d},
\]
the \emph{determinant} is $\det A = ad-bc$. We will build the determinant
properly in a later chapter (it measures signed area), but you can already
use it as a collapse-detector.

\begin{formalbox}
The $2\times 2$ matrix $A=\colvec{a&b\\c&d}$ is \textbf{nonsingular}
(invertible) iff $\det A = ad-bc \neq 0$, and \textbf{singular} iff
$\det A = 0$. When nonsingular, $A\vec x=\vec b$ has the unique solution
\[
  \vec x = A\inv \vec b,\qquad
  A\inv = \frac{1}{ad-bc}\colvec{d & -b\\ -c & a}.
\]
\end{formalbox}

Let us see \emph{why} $ad-bc$ is the right number, geometrically. The two
columns of $A$ are $\vec a_1=\colvec{a\\c}$ and $\vec a_2=\colvec{b\\d}$.
These two arrows are the two sides of a parallelogram. The
\emph{area} of that parallelogram is exactly $\lvert ad-bc\rvert$. Now:
\begin{itemize}
  \item If the columns point in different directions, the parallelogram has
  positive area, $ad-bc\neq0$: \emph{spread out}, nonsingular.
  \item If the columns are collinear, the parallelogram is squashed flat,
  area $0$, so $ad-bc=0$: \emph{collapsed}, singular.
\end{itemize}

\begin{figure}[H]
\centering
\begin{tikzpicture}[scale=0.85]
  \begin{scope}
    \node at (1.6,3.5) {\small $\det\neq0$: real area};
    \draw[->,gray] (-0.3,0)--(3.2,0); \draw[->,gray] (0,-0.3)--(0,3.2);
    \fill[vecblue!18] (0,0)--(2.2,0.6)--(2.9,2.6)--(0.7,2.0)--cycle;
    \draw[->,vecblue,very thick] (0,0)--(2.2,0.6) node[right]{\small $\vec a_1$};
    \draw[->,targetgreen!80!black,very thick] (0,0)--(0.7,2.0) node[above]{\small $\vec a_2$};
    \node at (1.45,1.3) {\small area $=|ad-bc|$};
  \end{scope}
  \begin{scope}[xshift=6cm]
    \node at (1.6,3.5) {\small $\det=0$: flat};
    \draw[->,gray] (-0.3,0)--(3.2,0); \draw[->,gray] (0,-0.3)--(0,3.2);
    \draw[gray!45,dashed] (0,0)--(3.0,2.4);
    \draw[->,vecblue,very thick] (0,0)--(1.5,1.2) node[below right]{\small $\vec a_1$};
    \draw[->,targetgreen!80!black,very thick] (0,0)--(2.5,2.0) node[above left]{\small $\vec a_2$};
    \node at (1.7,0.5) {\small area $=0$};
  \end{scope}
\end{tikzpicture}
\caption{The determinant is the signed area of the column parallelogram.
Nonzero area means the columns span the plane (one solution for every
$\vec b$); zero area means they have collapsed to a line (the singular
case).}
\label{ch06:detarea}
\end{figure}

Check our three examples against $ad-bc$:
\begin{itemize}
  \item Flagship \eqref{ch06:flagship}: $\det\colvec{2&1\\3&2}=2\cdot2-1\cdot3=1\neq0$.
  Nonsingular -- one solution. \checkmark
  \item No-solution \eqref{ch06:nosol}: $\det\colvec{1&2\\1&2}=1\cdot2-2\cdot1=0$.
  Singular -- and here $\vec b$ missed the line, so zero solutions.
  \item Infinite \eqref{ch06:infsol}: $\det\colvec{1&2\\2&4}=1\cdot4-2\cdot2=0$.
  Singular -- and here $\vec b$ sat on the line, so infinitely many.
\end{itemize}
The determinant alone separates ``nonsingular'' (always one solution) from
``singular'' (either zero or infinitely many, depending on $\vec b$). To
decide between those last two you must look at $\vec b$; the determinant
cannot, because it knows only about $A$.

\subsection{Deriving the inverse formula (so it is not magic)}

Where does $A\inv=\frac{1}{ad-bc}\colvec{d&-b\\-c&a}$ come from? Let us solve
$\colvec{a&b\\c&d}\colvec{x\\y}=\colvec{e\\f}$ by old-fashioned elimination,
assuming $ad-bc\neq0$. The equations are
\[
  ax+by=e,\qquad cx+dy=f.
\]
Multiply the first by $d$, the second by $b$, and subtract to kill $y$:
\[
  (ad-bc)x = de - bf \;\Longrightarrow\; x=\frac{de-bf}{ad-bc}.
\]
Multiply the first by $c$, the second by $a$, and subtract to kill $x$:
\[
  (bc-ad)y = ce - af \;\Longrightarrow\; y=\frac{af-ce}{ad-bc}.
\]
Stacking these,
\[
  \colvec{x\\y}=\frac{1}{ad-bc}\colvec{d & -b\\ -c & a}\colvec{e\\f},
\]
which is precisely $\vec x=A\inv\vec b$. The formula was never magic; it is
two lines of elimination with the division by $ad-bc$ done once at the end.
And now the danger of $ad-bc=0$ is glaring: it is the forbidden
\emph{division by zero}. Singular matrices have no inverse for the same
reason $\tfrac10$ has no value.

\begin{trapbox}
``$\det A=0$ means no solution'' is \emph{wrong}, and it is a trap students
fall into constantly. $\det A = 0$ means the unique-solution machine breaks;
the system then has \emph{either} no solution \emph{or} infinitely many. You
must inspect $\vec b$ (does it lie in the span of the columns?) to know
which. Nonsingular is a property of $A$ alone; the zero-versus-infinity split
needs $\vec b$.
\end{trapbox}

\begin{connectionbox}
The quantity $ad-bc$ is your first determinant, and determinants reach far
beyond linear algebra. In multivariable \emph{calculus}, the Jacobian
determinant is exactly this area factor: it tells you how a change of
variables stretches or squashes little area elements, which is why
$dx\,dy = \lvert\det J\rvert\,du\,dv$ appears in every change-of-variables
integral. A coordinate change is ``nonsingular'' (locally invertible)
precisely where its Jacobian determinant is nonzero -- the same
spread-out-versus-collapsed dichotomy you are learning here, now applied to
infinitesimal boxes.
\end{connectionbox}

\section{Up to three dimensions: planes instead of lines}

Everything generalizes cleanly to three unknowns; only the cast of
characters grows. An equation like $a x + b y + c z = d$ in three unknowns is
no longer a line -- it is a \emph{plane} in $\R^3$. A system of three such
equations asks for a point lying on all three planes at once.

\begin{intuitionbox}
In $\R^3$, each linear equation is a plane. The row picture of a
$3\times3$ system is \emph{three planes}, and a solution is a point common to
all three. Generically three planes meet at a single point -- the corner of
a room where two walls and the floor come together.
\end{intuitionbox}

\begin{figure}[H]
\centering
\begin{tikzpicture}[scale=1.5,line cap=round]
  % three "walls" meeting at a corner, schematic
  \coordinate (O) at (0,0,0);
  % floor
  \fill[targetgreen!18] (0,0,0)--(1.6,0,0)--(1.6,0,1.6)--(0,0,1.6)--cycle;
  % back wall
  \fill[vecblue!16] (0,0,0)--(1.6,0,0)--(1.6,1.6,0)--(0,1.6,0)--cycle;
  % side wall
  \fill[vecred!14] (0,0,0)--(0,0,1.6)--(0,1.6,1.6)--(0,1.6,0)--cycle;
  \draw[thick] (0,0,0)--(1.6,0,0);
  \draw[thick] (0,0,0)--(0,1.6,0);
  \draw[thick] (0,0,0)--(0,0,1.6);
  \filldraw[black] (0,0,0) circle (1.2pt);
  \node[below left] at (0,0,0) {\small solution};
  \node[targetgreen!60!black] at (1.1,0,1.1) {\small plane 3};
  \node[vecblue] at (1.2,1.1,0) {\small plane 1};
  \node[vecred!80!black] at (0,1.1,1.1) {\small plane 2};
\end{tikzpicture}
\caption{The row picture in $\R^3$: three planes meeting at one point, like
two walls and the floor meeting at the corner of a room. That corner is the
unique solution.}
\label{ch06:corner}
\end{figure}

The column picture also survives intact, and arguably gets \emph{more}
useful. A $3\times3$ system
\[
  A\vec x=\vec b,\qquad
  x\colvec{a_{11}\\a_{21}\\a_{31}} + y\colvec{a_{12}\\a_{22}\\a_{32}}
  + z\colvec{a_{13}\\a_{23}\\a_{33}} = \vec b,
\]
asks: can the target arrow $\vec b\in\R^3$ be built as a combination of the
three column arrows? If the three columns are \emph{spread out} -- not lying
in any common plane -- they span all of $\R^3$, and every $\vec b$ is
reachable in exactly one way. That is the nonsingular case.

\subsection{A worked $3\times3$ with a clean answer}

\begin{example}[Three planes, one corner]
\label{ch06:ex-3x3}
Solve
\[
  \begin{aligned}
    x + y + z &= 6,\\
    \phantom{x +{}} 2y + z &= 5,\\
    \phantom{x + y +{}} 2z &= 4.
  \end{aligned}
\]
This system is already in ``triangular'' shape, so we can read it from the
bottom up -- \emph{back-substitution}. The last equation gives $2z=4$, so
$z=2$. Plug into the middle: $2y+2=5$, so $y=\tfrac32$. Plug both into the
top: $x+\tfrac32+2=6$, so $x=\tfrac52$. The unique solution is
$(x,y,z)=\bigl(\tfrac52,\tfrac32,2\bigr)$.

The determinant of a triangular matrix is the product of its diagonal,
$1\cdot2\cdot2=4\neq0$, so the system is nonsingular -- three planes truly
crossing at the single corner $\bigl(\tfrac52,\tfrac32,2\bigr)$. Triangular
systems are the reward at the end of Gaussian elimination, which is the
subject of the next chapter; here we simply enjoy how easy they are.
\end{example}

\subsection{How a $3\times3$ system can fail}

In $\R^3$ the ways to fail are richer, but they still sort into ``no
solution'' and ``infinitely many,'' both flavors of \emph{singular}.

\textbf{No solution.} Three planes can fail to share a common point in
several shapes -- two of them parallel, or three forming a triangular prism
(each pair meets in a line, but no point is on all three, like the three
pages of an open book that refuse to share a single line). Algebraically,
elimination produces a contradiction such as $0=1$.

\textbf{Infinitely many.} If the three planes share a whole \emph{line}
(think of three pages of a book hinged on the same spine), every point on
that line solves the system: a line's worth of solutions. If they all
coincide as the \emph{same} plane, you get a plane's worth. Elimination
reveals this as a ``free variable'' you can set to anything.

\begin{figure}[H]
\centering
\begin{tikzpicture}[scale=1.0]
  % three planes sharing a common line (book spine), seen edge-on
  \draw[vecblue,very thick] (-1.5,0)--(1.5,0);
  \draw[vecred,very thick] (-1.4,-0.7)--(1.4,0.7);
  \draw[structurepurple,very thick] (-1.0,-1.3)--(1.0,1.3);
  \filldraw[black] (0,0) circle (2pt);
  \node[right] at (1.5,0) {\small \emph{spine}: shared line};
  \node at (0,-1.7) {\small three planes (edge-on) hinged on one line};
\end{tikzpicture}
\caption{An infinitely-many case in $\R^3$: three planes seen edge-on, all
hinged along a common line (a ``book spine''). Every point of that line
solves all three equations.}
\label{ch06:spine}
\end{figure}

\begin{focusbox}
The $0$/$1$/$\infty$ trichotomy is dimension-independent. In $\R^n$:
\begin{itemize}
  \item \textbf{Nonsingular} ($\det A\neq0$, columns span $\R^n$): exactly
  one solution for \emph{every} $\vec b$.
  \item \textbf{Singular} ($\det A=0$, columns trapped in a smaller flat):
  \emph{either} no solution (when $\vec b$ escapes that flat) \emph{or}
  infinitely many (when $\vec b$ lies inside it).
\end{itemize}
Same story as $2\times2$, with ``line'' upgraded to ``plane'' and
``parallelogram'' to ``parallelepiped'' whose volume is $\det A$.
\end{focusbox}

\section{The homogeneous shadow: $A\vec x=\vec 0$}

There is a special right-hand side that controls all the others: the zero
vector. The system $A\vec x=\vec 0$ is called \emph{homogeneous}. It always
has at least the boring solution $\vec x=\vec 0$ (do nothing, reach the
origin). The real question is whether it has \emph{others}.

\begin{theorem}[Singular $\Leftrightarrow$ nontrivial null solution]
\label{ch06:homog}
A square matrix $A$ is singular if and only if $A\vec x=\vec 0$ has a
solution $\vec x\neq\vec 0$.
\end{theorem}

\begin{proof}[Sketch]
If $A$ is nonsingular it has an inverse, so $\vec x=A\inv\vec 0=\vec 0$ is
the \emph{only} solution -- no nontrivial one exists. Conversely, if the
columns of $A$ are collapsed (singular), some nontrivial combination of them
equals $\vec 0$; the weights of that combination form an $\vec x\neq\vec0$
with $A\vec x=\vec0$. For our flagship $A=\colvec{2&1\\3&2}$ (nonsingular),
only $\vec0$ works. For the collapsed $\colvec{1&2\\2&4}$, the combination
$2\colvec{1\\2}-1\colvec{2\\4}=\vec0$ exposes $\vec x=\colvec{2\\-1}\neq\vec0$.
\end{proof}

This is the bridge between the row and column pictures and the rest of the
book. The nontrivial solutions of $A\vec x=\vec 0$ form the \emph{null
space} of $A$; a nonsingular matrix has a trivial null space, a singular one
does not. And here is the elegant payoff connecting back to
Proposition~\ref{ch06:zoi}:

\begin{aestheticbox}
Every solution of $A\vec x=\vec b$ has the shape
\[
  \vec x = \underbrace{\vec x_{\text{particular}}}_{\text{one solution}}
         + \underbrace{\vec x_{\text{null}}}_{\text{anything killed by }A}.
\]
The particular part gets you to $\vec b$; the null-space part is the freedom
to wander without changing where you land. If the null space is just
$\{\vec0\}$ (nonsingular), there is no freedom -- one solution. If the null
space is a line or a plane (singular), that same flat shape is copied onto
the solution set, giving infinitely many. \emph{The geometry of the
solutions is the geometry of the null space, shifted.} This one decomposition
quietly organizes the chapters to come.
\end{aestheticbox}

Figure~\ref{ch06:shift} draws this for a singular system whose null space is
a line: the full solution set is that same line, picked up and slid over so
it passes through one particular solution $\vec x_{\text{particular}}$.

\begin{figure}[H]
\centering
\begin{tikzpicture}[scale=0.9]
  \draw[->,gray] (-2.6,0)--(2.8,0) node[right]{};
  \draw[->,gray] (0,-1.6)--(0,2.9) node[above]{};
  % null space: a line through the origin, direction (1,1)
  \draw[structurepurple,very thick] (-1.8,-1.8)--(1.8,1.8);
  \node[structurepurple] at (1.9,1.45) {\small null space};
  \filldraw[black] (0,0) circle (1.6pt) node[below right]{\small $\vec 0$};
  % particular solution
  \draw[->,vecblue,very thick] (0,0)--(-1.2,1.3) node[midway,above left]{\small $\vec x_{\text{part}}$};
  \filldraw[vecblue] (-1.2,1.3) circle (1.6pt);
  % solution set = null space shifted to pass through x_part: direction (1,1)
  \draw[vecred,very thick,dashed] (-2.5,0.0)--(0.6,3.1);
  \node[vecred] at (0.95,2.85) {\small solution set};
\end{tikzpicture}
\caption{Solution set $=$ null space, shifted. The solid purple line is the
null space (all $\vec x$ with $A\vec x=\vec0$). The dashed red line is the
full solution set of $A\vec x=\vec b$: the \emph{same direction}, slid over to
pass through one particular solution $\vec x_{\text{part}}$. Empty null space
(a single dot at $\vec 0$) would shift to a single solution.}
\label{ch06:shift}
\end{figure}

\begin{connectionbox}
This ``particular plus homogeneous'' shape is not special to matrices -- it
is the signature of \emph{any} linear operator. In a first course on
\emph{differential equations} you meet exactly the same theorem: the general
solution of $y'' + y = \cos t$ is one particular solution (say
$y_p=\tfrac{t}{2}\sin t$) plus the general solution of the homogeneous
equation $y''+y=0$ (here $c_1\cos t + c_2\sin t$). The homogeneous solutions
play the role of the null space, and the particular solution plays the role
of $\vec x_{\text{part}}$. The map ``apply the differential operator
$L=\tfrac{d^2}{dt^2}+1$'' is linear in precisely the sense
$L(\alpha f+\beta g)=\alpha L f+\beta Lg$, so the proof is word-for-word the
one we just gave for $A\vec x=\vec b$. Once you see that solution sets of all
linear problems are ``a null space, shifted,'' the two subjects fuse.
\end{connectionbox}

\section{A modern application: when ``solve $A\vec x=\vec b$'' is the whole job}

\begin{applicationbox}
\textbf{Networks, circuits, and the structure under deep learning.}
Linear systems are not a warm-up; they are the load-bearing computation in
enormous parts of technology.

\emph{Circuits and networks.} Kirchhoff's laws for an electrical network --
``current in equals current out at every node, voltage drops around every
loop sum to zero'' -- are a linear system $A\vec x=\vec b$ whose unknowns are
the currents. The \emph{same} matrix shape (a graph's incidence matrix)
governs water in pipes, traffic on roads, and forces in a bridge truss.
Engineers solving for ``will this bridge stand?'' are, underneath, asking
whether a structural stiffness matrix is nonsingular and what $A\inv\vec b$
says about deflections.

\emph{Computer graphics.} Every time a 3D scene is rendered, the engine
solves small linear systems to find where a ray pierces a triangle, or to
blend bone transforms for an animated character. The ``does the ray hit?''
test is exactly a $3\times3$ singular-versus-nonsingular check: a singular
system means the ray is parallel to the surface and grazes by.

\emph{Machine learning.} Fitting a linear model -- the regression that
underlies a shocking amount of applied ML -- is set up as $A\vec x=\vec b$,
where rows are data examples and $\vec x$ is the vector of weights. Usually
there are far more equations than unknowns and \emph{no exact solution
exists} (the $\vec b$ of real data does not lie in the column space). The
fix, \emph{least squares}, finds the $\vec x$ that gets closest -- it asks
``what is the nearest reachable point to $\vec b$ in the span of the
columns?'' That question is the column picture of this chapter, asked when
the target is just out of reach. Even the giant nonlinear models you hear
about are trained by repeatedly solving (or approximately solving) linear
systems for the update direction. The trichotomy you learned today -- spread
out vs.\ collapsed columns -- reappears as ``is this problem
well-conditioned, or has the data secretly collapsed onto a lower-dimensional
flat?'' That collapse, called multicollinearity, is one of the most common
ways a real-world model goes haywire.
\end{applicationbox}

Figure~\ref{ch06:leastsq} shows the least-squares idea in one picture. The
reachable targets -- all combinations of the columns -- fill a plane (the
column space). When the data vector $\vec b$ floats \emph{off} that plane, no
exact solution exists, so we drop a perpendicular and solve for the nearest
reachable point $\hat{\vec b}$ instead. That ``drop a perpendicular onto the
column space'' is the whole of linear regression in one move; we make it
precise when we reach projections.

\begin{figure}[H]
\centering
\begin{tikzpicture}[scale=1.0,line cap=round]
  % the column-space plane, drawn as a slanted parallelogram
  \fill[targetgreen!14] (-2.2,-0.7)--(2.2,-0.7)--(3.0,0.9)--(-1.4,0.9)--cycle;
  \draw[targetgreen!60!black] (-2.2,-0.7)--(2.2,-0.7)--(3.0,0.9)--(-1.4,0.9)--cycle;
  \node[targetgreen!55!black] at (2.0,-0.45) {\small column space (reachable $\vec b$)};
  \coordinate (O) at (0,0);
  \coordinate (B) at (1.1,2.5);   % b off the plane
  \coordinate (P) at (1.1,0.35);  % foot of perpendicular (projection)
  \draw[->,vecred,ultra thick] (O)--(B) node[above right]{\small $\vec b$ (data, off-plane)};
  \draw[->,vecblue,very thick] (O)--(P) node[below right]{\small $\hat{\vec b}=A\hat{\vec x}$};
  \draw[trapred,thick,dashed] (B)--(P);
  \draw[trapred] ($(P)+(-0.16,0)$)--($(P)+(-0.16,0.16)$)--($(P)+(0,0.16)$);
  \node[trapred!85!black] at (1.95,1.5) {\small residual $\vec b-\hat{\vec b}$};
  \filldraw[black] (O) circle (1.3pt);
\end{tikzpicture}
\caption{Least squares as the column picture with the target out of reach.
The reachable vectors form a flat (the column space). When $\vec b$ lies off
it, we solve $A\vec x=\hat{\vec b}$ for the nearest point $\hat{\vec b}$, the
perpendicular foot. The dashed residual is what regression makes as small as
possible.}
\label{ch06:leastsq}
\end{figure}

\section{Building intuition that lasts}

\begin{walkbox}
Take a walk and carry exactly this image. You are handed two arrows on the
ground, $\vec a_1$ and $\vec a_2$, and a flag planted at point $\vec b$. Your
job: reach the flag by walking some amount along $\vec a_1$ and some amount
along $\vec a_2$ (backwards is allowed). Ask yourself:
\begin{itemize}
  \item If the two arrows point in genuinely different directions, can you
  always reach \emph{any} flag? In how many ways?
  \item If the two arrows point the \emph{same} way, which flags can you
  reach, and which are forever out of reach?
\end{itemize}
When you can answer these without paper, you understand the column picture --
and therefore the existence-and-uniqueness theory of linear systems -- better
than most people who have passed the course.
\end{walkbox}

\begin{remark}[Two pictures, one habit]
The professional habit this chapter is trying to install: whenever you see
$A\vec x=\vec b$, \emph{flicker} between the two pictures. Need to know if a
solution \emph{exists}? Think columns -- is $\vec b$ in their span? Need to
\emph{compute} the solution? Think rows -- eliminate. Need to know if the
solution is \emph{unique}? Ask whether the columns are spread out
($\det\neq0$) or collapsed. The fluent reader does this flicker
unconsciously, many times a page, for the rest of the book.
\end{remark}

\section{Summary}

\begin{itemize}
  \item A linear system $A\vec x=\vec b$ has \textbf{two geometric readings}.
  The \emph{row picture} sees each equation as a line (in $\R^2$) or plane
  (in $\R^3$); a solution is a common intersection point. The \emph{column
  picture} sees $A\vec x$ as a linear combination of the columns; a solution
  is a recipe of weights reaching $\vec b$.
  \item Compute $A\vec x$ two ways: \emph{rows dotted with $\vec x$}, or
  \emph{a combination of the columns}. Both are the same product.
  \item Exactly three outcomes are possible: \textbf{one}, \textbf{none}, or
  \textbf{infinitely many} solutions. Two distinct solutions force a whole
  line of them (Proposition~\ref{ch06:zoi}).
  \item \textbf{Nonsingular} ($\det A\neq0$): columns spread out / span the
  space; one solution for every $\vec b$; $\vec x=A\inv\vec b$.
  \textbf{Singular} ($\det A=0$): columns collapsed onto a smaller flat;
  then \emph{no} solution if $\vec b$ escapes that flat, \emph{infinitely
  many} if $\vec b$ lies inside it.
  \item The determinant $ad-bc$ (in $2\times2$) is the signed area of the
  column parallelogram: zero area = collapse = singular.
  \item The homogeneous system $A\vec x=\vec0$ measures the freedom: its
  nonzero solutions (the null space) are exactly the wiggle room that turns
  one solution into infinitely many.
\end{itemize}

\begin{marginnotebox}
\textbf{Cost and intuition recap.} The $2\times2$ inverse formula costs a
handful of multiplications, and Cramer-style formulas exist in every
dimension -- but their cost explodes (an $n\times n$ determinant by the naive
expansion costs on the order of $n!$ operations, hopeless past $n\approx 10$).
The practical route is \emph{elimination}, which solves $A\vec x=\vec b$ in
about $\tfrac13 n^3$ operations -- the workhorse of the next chapter. So carry
two things forward: the \emph{intuition} (rows are lines/planes, columns are
arrows, singular means collapsed) and the awareness that \emph{how} we
compute the answer, not just whether it exists, will matter enormously once
$n$ is large.
\end{marginnotebox}

\begin{historybox}
The systematic solution of linear systems is one of the oldest mathematical practices in human history.
Long before European algebraists began manipulating symbols, Chinese civil servants in the Han Dynasty
(compiled around 150~BCE by Zhang Cang and Geng Shoushan) recorded the method in \emph{Jiuzhang Suanshu}
(The Nine Chapters on the Mathematical Art). In Chapter~8, titled \emph{Fangcheng} (``The Rectangular Array''),
the authors solved a system of three linear equations in three agricultural unknowns (bundles of good, medium,
and poor grain) by laying out the coefficients on a counting board using bamboo rods.

The Han mathematicians arranged equations into columns on the board and used red bamboo rods for positive
coefficients and black rods for negative coefficients -- the earliest recorded use of negative numbers in
the history of mathematics, centuries before negative quantities were accepted in Europe. By repeatedly
subtracting multiples of one column from another, they reduced the array to upper-triangular form and read off
the solution by back-substitution.

For centuries in the West, linear systems were treated as pure algebraic drudgery -- variable substitution
and cancellation. When geometry was introduced, it was almost always the row picture: two lines
meeting at a point, or three planes meeting at a corner. But the row picture has a fatal limitation:
it collapses in higher dimensions. While human intuition can picture three planes intersecting in 3D space,
nobody can visualize four hyperplanes intersecting in four-dimensional space $\R^4$, let alone fifty hyperplanes
in $\R^{50}$.

The revolutionary insight championed by modern educators like Gilbert Strang is to elevate the \textbf{column picture}:
\[
  x_1\vec{a}_1 + x_2\vec{a}_2 + \cdots + x_n\vec{a}_n = \vec{b}.
\]
In the column picture, the geometry does not change as the dimension $n$ increases! In $\R^2$, $\R^3$, or $\R^{1000}$,
the column picture always asks the exact same intuitive question: can you scale the given arrows $\vec{a}_i$ and
lay them tip-to-tail to reach the target arrow $\vec{b}$? Solvability is simply reachability in the column space.

Framing a linear system as $\mat{A}\vec{x} = \vec{b}$ connects algebraic equations directly to
the geometry of vectors and spans, giving us a definitive taxonomy of solutions: exactly one (columns span space
without redundancy), none (target is outside the column span), or infinitely many (redundancy in the columns creates
a nullspace of non-trivial solutions to $\mat{A}\vec{x}=\vec{0}$). Yet drawing intersecting hyperplanes or column
vectors provides conceptual clarity, not an automated computational tool for systems with hundreds of variables.
Manual variable substitution quickly becomes unmanageable, while naive determinant expansions like Cramer's rule
explode factorially in complexity ($O(n!)$ operations), quickly overwhelming any computer. Transforming systems
reliably requires a mechanized, scalable, polynomial-time algorithm. That engine arrives in Chapter~\ref{ch07:two-lines}
with Gaussian elimination, which systematically reduces any linear system to triangular form in $O(n^3)$ steps.
\end{historybox}

\section{Exercises}

\begin{exercise}[Both pictures, by hand]
\label{ch06:ex-both}
For the system $\;3x+y=9,\ \ x+2y=8$:
(a) draw the row picture (two lines) and read off the crossing point;
(b) write it as $x\colvec{3\\1}+y\colvec{1\\2}=\colvec{9\\8}$ and find the
weights; (c) confirm both answers agree.
\end{exercise}

\begin{exercise}[The determinant decides]
\label{ch06:ex-det}
Compute $\det A$ for each matrix and state, \emph{without solving}, whether
$A\vec x=\vec b$ is guaranteed to have a unique solution for every $\vec b$:
\[
  \colvec{4 & 2\\ 1 & 3},\qquad
  \colvec{2 & 6\\ 1 & 3},\qquad
  \colvec{0 & 5\\ -2 & 0}.
\]
For any singular one, find a $\vec b$ giving no solution and a $\vec b$
giving infinitely many.
\end{exercise}

\begin{exercise}[Make it fail two ways]
\label{ch06:ex-fail}
Start from $\;x+3y=4$. Add a \emph{second} equation so that the system has
(a) no solution; (b) infinitely many solutions; (c) exactly one solution.
For each, name the geometric relationship between the two lines.
\end{exercise}

\begin{exercise}[Column reach]
\label{ch06:ex-reach}
The columns of $A$ are $\vec a_1=\colvec{1\\2}$ and $\vec a_2=\colvec{2\\4}$.
(a) Describe, as a set, all vectors $\vec b$ for which $A\vec x=\vec b$ has a
solution. (b) For $\vec b=\colvec{3\\6}$, give \emph{two} different solutions
$\vec x$, and then describe \emph{all} of them in the form
$\vec x_{\text{particular}}+t\,\vec x_{\text{null}}$.
\end{exercise}

\begin{exercise}[Back-substitution in $\R^3$]
\label{ch06:ex-back}
Solve, top-triangular style,
\[
  2x + y - z = 3,\qquad 3y + 2z = 4,\qquad 4z = 8,
\]
and verify your point satisfies all three equations. What is $\det$ of the
coefficient matrix, and why does that guarantee uniqueness?
\end{exercise}

\begin{exercise}[Three planes, no corner]
\label{ch06:ex-prism}
Consider $\;x+y+z=1,\ \ x+y+z=2,\ \ x-y=0$. Show by inspection that the
system has no solution, and explain in words what the three planes are doing
(which two are parallel?).
\end{exercise}

\begin{exercise}[Derive your own inverse]
\label{ch06:ex-inverse}
Without quoting the formula, solve $\;ax+by=e,\ cx+dy=f$ by elimination to
re-derive $x=\frac{de-bf}{ad-bc}$ and $y=\frac{af-ce}{ad-bc}$. At which step
does the assumption $ad-bc\neq0$ become essential, and what goes wrong if it
fails?
\end{exercise}

\begin{exercise}[Zero, one, or infinity -- a proof]
\label{ch06:ex-proof}
Suppose $A\vec p=\vec b$ and $A\vec q=\vec b$ with $\vec p\neq\vec q$. Prove
that $\vec p + t(\vec q-\vec p)$ solves $A\vec x=\vec b$ for \emph{every} real
$t$, citing the two linearity rules you use. Conclude that a linear system
cannot have exactly two solutions.
\end{exercise}

\begin{exercise}[The homogeneous shadow]
\label{ch06:ex-homog}
For $A=\colvec{1 & 2 & 3\\ 2 & 4 & 6\\ 1 & 1 & 1}$:
(a) notice that row $2$ is exactly twice row $1$, so $A$ cannot have full
rank; use this to argue $A$ is singular. (b) Now find a nonzero $\vec x$ with
$A\vec x=\vec0$. (Hint: a null vector records a dependence among the
\emph{columns}; look for weights $x_1,x_2,x_3$, not all zero, with
$x_1\colvec{1\\2\\1}+x_2\colvec{2\\4\\1}+x_3\colvec{3\\6\\1}=\vec0$. You can
check that $\vec x=\colvec{1\\-2\\1}$ works.) (c) What does the existence of
this $\vec x$ tell you about the number of solutions of $A\vec x=\vec b$
\emph{when a solution exists at all}?
\end{exercise}

\begin{exercise}[Challenge: a parametric crossing]
\label{ch06:ex-param}
Consider $\;x+y=1,\ \ x+ky=2$, where $k$ is a real parameter.
(a) For which value(s) of $k$ is the coefficient matrix singular?
(b) For that $k$, does the system have no solution or infinitely many?
(c) For all other $k$, solve for $(x,y)$ in terms of $k$ and describe what
happens to the solution geometrically as $k\to1$. (Watch the crossing point
fly off to infinity as the lines become parallel.)
\end{exercise}
